% Section 13 — Conclusion
\section{Conclusion}
\label{sec:conclusion}

AlphaGo Zero made search unnecessary for a strong move in one game. EurekaZero asks for the same
property in games whose boards are new problems: a dataset that needs a learning procedure, an
optimization history that needs its next experiment, a specification that needs a part. Under a
protocol that commits one candidate before anything is executed, scored or revised, the proposal is
the move, the verifier is the referee, and the best known solution is the opponent.

The evidence this draft describes would establish three things: that one interface can play
innovation games with different output roles; that its first proposal beats the strongest
non-search control and is worth tens of evaluations of search; and that some of its winning moves are
new in a sense that survives intervention and replication. The gap ledger (\cref{app:gaps}) is the
list of experiments between this draft and that paper. The program's own record (\cref{app:history})
shows why the list matters: it has already measured continued training on its own experience raising
regression alerts rather than improvements. If the confirmation runs reproduce that failure, this paper will report it.

If they do not, innovation that today needs a search for every problem would need it only once per
family, and the first proposal would arrive while the problem is still being described.
